% Inside a participant controller, and the one distinction this volume already made.
\begin{tikzpicture}[font=\sffamily\scriptsize, x=1cm, y=1cm]

\node[chlabel, anchor=west] at (-6.8,3.4) {inside a participant controller, and the part of it chapter 11 already reinvented};

% ================================================================ the blocks
\node[hw, text width=24mm, minimum height=7mm] (port) at (-5.4,2.5) {two ports, minimum};
\node[hw, text width=26mm, minimum height=7mm] (fwd) at (-2.2,2.5)
  {auto-forwarder and loopback};
\node[hw, text width=26mm, minimum height=7mm] (proc) at (1.2,2.5) {the processing unit};
\draw[flow] (port) -- (fwd);
\draw[flow] (fwd) -- (proc);
\node[chlabel, anchor=west, text width=32mm] at (2.9,2.5)
  {unconnected ports close their loops automatically};

\node[kern, text width=26mm, minimum height=7mm] (fmmu) at (-4.4,1.1)
  {address mapping, bit wise};
\node[kern, text width=26mm, minimum height=7mm] (sm) at (-1.0,1.1) {sync managers};
\node[kern, text width=26mm, minimum height=7mm] (ram) at (2.4,1.1) {process data memory};
\draw[flow] (proc.south) -- (1.2,1.7) -- (-4.4,1.7) -- (fmmu.north);
\draw[flow] (-1.0,1.7) -- (sm.north);
\draw[flow] (2.4,1.7) -- (ram.north);

\node[hw, text width=26mm, minimum height=7mm] (dc) at (-4.4,-0.2) {distributed clocks};
\node[hw, text width=26mm, minimum height=7mm] (pdi) at (-1.0,-0.2)
  {the process data interface};
\node[fw, text width=26mm, minimum height=7mm] (mcu) at (2.4,-0.2) {a host processor};
\draw[flow] (pdi) -- (mcu);
\draw[flow] (sm) -- (pdi);
\node[chlabel, anchor=west, text width=32mm] at (4.0,-0.2)
  {ceiling: about 12.5 megabytes per second};

% ================================================================ the two modes
\node[regcell, fill=usrcol, minimum width=54mm, text width=52mm, inner sep=3pt,
      anchor=north west, align=left] at (-6.8,-1.4)
  {\textbf{buffered mode}: three buffers, cyclic data, and the consumer always
   receives the latest \textbf{consistent} buffer. A late reader gets fresh data
   and never a torn one.\\[2pt]
   \textbf{This is chapter 11's state frame, exactly.}};

\node[regcell, fill=kercol, minimum width=54mm, text width=52mm, inner sep=3pt,
      anchor=north west, align=left] at (0.2,-1.4)
  {\textbf{mailbox mode}: a handshake, and nothing is lost.\\[2pt]
   \textbf{This is chapter 11's command and configuration path}, where a
   dropped message is a fault rather than a stale value.};

\node[umlnote, text width=114mm, anchor=north west] at (-6.8,-4.0)
  {Chapter 11 reached the same split from first principles over a different bus:
   state that may be superseded and must never be torn, against commands that
   must not be lost. Finding the same distinction implemented in silicon, by
   people solving the same problem with a hardware budget, is the strongest
   independent confirmation that design has had. It is also the reason this
   chapter is worth reading even for somebody who will never buy the chip.};
\end{tikzpicture}
